\section{Experimental Setup}
\label{sec:setup}

% Evidence: canonical/data.json#protocol; experiments/configs/phase2|3|5;
% experiments/results/phase4/comparison.json#protocol; train.py transforms/optimizer;
% phase2 yaml environment; paper/conclusion.md §2 (hardware).
% Missing at write time: research_os/paper/outline.md, claim_ledger.md, citation_plan.md.

\subsection{Dataset and preprocessing}
\label{sec:setup-data}

We use \textbf{CIFAR-10} with the official $50\mathrm{k}/10\mathrm{k}$ train/test split.
Training images use random crop with padding~$4$ ($\texttt{RandomCrop}(32,\texttt{padding}{=}4)$),
random horizontal flip, conversion to tensors, and per-channel normalization with
mean $(0.4914, 0.4822, 0.4465)$ and std $(0.2023, 0.1994, 0.2010)$
as implemented in the training script.
Test images use tensor conversion and the same normalization without spatial augmentation.

\subsection{Optimization and schedule}
\label{sec:setup-optim}

Unless stated otherwise, all reported models share the following frozen recipe
(canonical protocol; Phase~2--5 configs):
\begin{itemize}
  \item \textbf{Optimizer:} SGD with momentum $0.9$ and weight decay $10^{-4}$.
  \item \textbf{Learning rate:} $0.1$, decayed by $10\times$ at iteration milestones
        $32{,}000$ and $48{,}000$ (\texttt{MultiStepLR}).
  \item \textbf{Mini-batch size:} $128$.
  \item \textbf{Training length:} $160$ epochs.
\end{itemize}
Configs prescribe the same recipe for Plain-20, ResNet-20-A, PreAct-ResNet-20-A, and Ablations~A/B;
Phase~4 comparisons are read-only aggregations of existing runs (no re-training).
Runtime identity of every transform/optimizer call across Phase~3/5 logs is not exhaustively
re-audited here
\texttt{[NEEDS\_VERIFICATION]}.

\subsection{Evaluation protocol}
\label{sec:setup-eval}

The primary metric is \textbf{best test error} ($\%$).
The implementation evaluates on the test set every $4$ epochs and selects
\texttt{best.pt} by minimum test error
(matching the reference training script's checkpoint rule).
We report the best test errors
$\ErrPlain\%$ (Plain-20),
$\ErrResNet\%$ (ResNet-20-A),
$\ErrPreAct\%$ (PreAct-ResNet-20-A), and
$\ErrAblationA\%$ (Ablation~A)
from the paper canonical snapshot
(\texttt{research\_os/paper/canonical/data.json}),
which is built from
\texttt{experiments/manifests/} and \texttt{experiments/results/}.
Ablation~B is recorded as a failed run (\AblationBStatus) and is \emph{not} assigned a
comparable best-test-error number.

\subsection{Randomness and hardware}
\label{sec:setup-hw}

\paragraph{Single-seed limitation.}
No random seed is fixed in the recorded protocol ($\texttt{seed}{=}\texttt{null}$).
Each configuration is a single completed training trajectory unless noted as failed.
Therefore we do not report mean$\pm$std over seeds, and small deltas
(e.g., Ablation~A vs.\ ResNet-20-A) must not be read as statistically tested improvements.

\paragraph{Hardware and software.}
Recorded notes for the frozen baseline protocol include an
NVIDIA GeForce RTX~3050 Laptop GPU (4\,GB) and PyTorch $2.0.1{+}\mathrm{cu117}$
under conda environment \texttt{deeplearn}
(Phase~2 config / Phase~4 protocol notes; also summarized in \texttt{paper/conclusion.md}).
We do \emph{not} claim that every Phase~3/5 canonical run used identical driver/CUDA build strings;
an earlier environment note also mentions a uv/cu118 path that may differ from the conda pin
\texttt{[NEEDS\_VERIFICATION]}.

\subsection{Scope boundaries for interpretation}
\label{sec:setup-scope}

All setup claims in this section are limited to CIFAR-10, depth~20, Option~A residual models,
the frozen post-activation training recipe above, and single-run evaluation.
We do not claim a full multi-depth Table~6 reproduction suite, ImageNet results,
or recipe retuning specialized to PreAct.
